\documentclass[a4paper,12pt]{article}
\usepackage{amsmath}
\usepackage{bm}
\usepackage[mathscr]{euscript}
\usepackage{framed}
\usepackage{graphicx}

\begin{document}

\title{Programming Exam Documentation}

\author{Marcel Dürr \and Enes Witwit}

\pagestyle{empty}


%%%%%%%%%%%%%%%%%%%%%%%%%%%%%%%%%%%%%%%%%%%%%%%%%%%%%%%%%%%%%%%%%%


% Titelblatt der Arbeit
\begin{titlepage}

%Logo
\begin{center} 
{\Huge Numerical methods for partial differential equations}
\vspace*{1cm}

{\Large Programming Exam}
\vspace*{1.5cm}

by \\[0.5cm]

{\Large Enes Witwit and Marcel Duerr}
\vspace*{1.5cm}

Examined by  \\[0.5cm]
{\Large Prof. Dr. Kanschat}
\vspace*{3cm}

Faculty for math and computer science\\[0.5cm]
University of Heidelberg \\[0.8cm]
\includegraphics[scale=0.2]{hd-logo.jpg} 
\end{center}
\end{titlepage}

\section{Introduction}
The programming exam states the following task.
\begin{framed}
Write a finite element code in your preferred programming language solving the boundary value problem:
\begin{align}
\mbox{Find } u\in C^2(\Omega)\mbox{ :}-\Delta u+u &= \cos(\pi x)\cos(\pi y) &&\mbox{in } \Omega \\
\partial _n u &= 0 &&\mbox{on } \partial\Omega.
\end{align}
Here, $\Omega = [0,1]^2$.
\end{framed}
First of all, we want to examine the given BVP. In order to produce meaningful results, we have to proof that a unique solution exists. Furthermore, this document will explain some of the specific formulars used for computation. These formulars depend heavily on the structure of the problem and can not be generalized.

\section{Weak Formulation}
In order to apply some of the results from the lecture, we need to derive the weak formulation of the given problem 
\begin{align}
\mbox{Find } u\in C^2(\Omega)\mbox{ :}-\Delta u+u &= \cos(\pi x)\cos(\pi y) &&\mbox{in } \Omega \\
\partial _n u &= 0 &&\mbox{on } \partial\Omega.
\end{align}
Multiplying with an arbitary $v\in C^2(\Omega)$ and integrating over $\Omega$ gives us
\[-\int _\Omega \Delta uv \,\mbox{d} \bm{x} + \int _\Omega uv \,\mbox{d} \bm{x} = \int _ \Omega fv\,\mbox{d} \bm{x}\]
where $f=\cos(\pi x)\cos(\pi y)$ and $\bm{x}=(x,y)$. Using Green's first formula and (2) we can obtain the weak formulation
\[\int _\Omega \nabla u\nabla v \,\mbox{d} \bm{x} + \int _\Omega uv \,\mbox{d} \bm{x} = \int _ \Omega fv\,\mbox{d} \bm{x}\]
From now on, we will denote the left hand side of the equation by $a(u,v)$ and the right hand side by $F(v)$. Thus, we obtain the weak formulation
\begin{equation}
\mbox{Find }u\in H^1(\Omega) \mbox{ :} \quad a(u,v)=F(v)\quad \forall v \in H^1(\Omega)
\end{equation}

\section{Existence and Uniqueness of a Solution}

We can already see, that our bilinear form $a(\cdotp ,\cdot)$ is the inner product associated with the norm on our function space $H^1(\Omega)$. We want to use the Riesz representation theorem to prove existence and uniqueness of a solution. In order to do so, it remains to show that our functional $F(\cdot)$ is linear and bounded. Linearity follows from the properties of integration. Using Hoelder's inequality, we show that
\begin{align*}
F(u)&=\|fu\|_{L^1(\Omega)}\\
	&\stackrel{\tiny\mbox{Hld.}}{\le} \|f\|_{L^2(\Omega)}\|u\|_{L^2(\Omega)}\\
	&\le \|1\|_{L^2(\Omega)}\ \big(\|u\|_{L^2(\Omega)}+\|u\|_{L^2(\Omega)}\big)\\
	&\le c\ \|u\|_{H^1(\Omega)}\mbox{,}
\end{align*}
where c depends on our domain $\Omega$. For our case, we have $\Omega=[0,1]^2$, in particular this means that $\Omega$ is bounded and our constant $c$ is finite. Therefore, $F(\cdot)$ is a bounded, linear functional and we can apply the Riesz representation theorem.

\section{Finding the Analytical Solution}

Now that we know that a unique solution exists, we want to actually compute it. We will use the ansatz $u=C\cos(\pi x)\cos(\pi y)$, with its gradient $\Delta u=2\pi^2 u$. Inserting in (1) gives us
\begin{equation}
-2C\pi^2 \cos(\pi x)\cos(\pi y) + C\cos(\pi x)\cos(\pi y)=\cos(\pi x)\cos(\pi y)
\end{equation}
\begin{align}
&\Rightarrow - 2C\pi^2 + C =1\\
&\Leftrightarrow C = \frac{1}{1-2\pi^2}
\end{align}
This leaves us with the solution $u=\frac{1}{1-2\pi^2}\cos(\pi x)\cos(\pi y)$.

\section{Integrals and Transformations}

The transformations we have to use for our FEM-code are very simple, since our mesh will be generated uniformly. Ultimately, we will only have to scale and displace our reference cell. We decided, that we will hard-code this property, since it will heavily simplifiy the computation of the local stiffness-matrices and the right-hand-side of the linear system, as we will see in brief.
\\ \\
Let $T$ be a cell of the mesh with vertices $v_1,v_2,v_3,v_4$, where $v_1$ is the vertex on the bottom left and h is the length of every edge. $\hat T = [0,1]^2$ shall be our reference cell. Then, our transformation $F$ will look like 
\[F \smash : \  \hat T \rightarrow T, \ \ \bm{\hat x}=\begin{pmatrix}\hat x \\ \hat y \end{pmatrix} \mapsto \begin{pmatrix}h \ 0 \\ 0 \ h\end{pmatrix} \begin{pmatrix}\hat x \\ \hat y \end{pmatrix} + v_1 := \begin{pmatrix}x\\y\end{pmatrix} := \bm{x}\]
We immediately can see, that the jacobian $J$ is given by
\begin{equation}
J(\bm{\hat x})=J=\begin{pmatrix}h \ 0\\0 \ h\end{pmatrix}\mbox{,}
\end{equation}
in particular, it is independet of $\bm{\hat x}$. Consequently, we also get
\begin{equation}
|\mbox{det}(J(\bm{\hat x}))|=|\mbox{det}(J)|=h^2\mbox{and} \ J^{-1}=h^{-1}\begin{pmatrix}1 \ 0 \\ 0 \ 1\end{pmatrix}
\end{equation}
Using the results from the lecture on the transformation of our bilinear form $a(\cdotp, \cdot)$ and using (7) and (8), we can compute the stiffness matrix in the following way
\[a^{(T)}_{ij} = \int_{\hat T}\nabla p_i \,\cdotp \nabla p_j\,\mbox{d}\bm{\hat x} + h^2\int_{\hat T}p_i p_j\,\mbox{d}\bm{\hat x}\mbox{,}\]
where $p_i$ and $p_j$ are our shape functions.\\ \\
\textit{Remark:} Note that the $\nabla$ in our formula refers to the gradient with respect to $\bm{\hat x}=\begin{pmatrix} \hat x \\ \hat y\end{pmatrix}$.
\textit{Remark:} Note that for affine transformations, the local stiffness matrices do not depend on the cell $T$, for which we want to compute it. Therefore, every local matrix looks the same and we only have to compute it once.
\\ \\
For the right-hand side, we have to look at the substitution formula for higher dimension integrals. But first of all, we need do split the intagral up. Let $\phi_i$ be a basis function of our domain $\Omega$, $\mathscr{T}_h$ our meshcells and $\mathscr{T}_i$ the mesh cells, on which $\phi_i \neq 0$. Then we get
\begin{align*}
\int_{\Omega}f(\bm{x})\phi_i(\bm{x})\ \mbox{d}\bm{x}&=\sum_{T\in\mathscr{T}_h} \int_Tf(\bm{x})\phi_i(\bm{x})\ \mbox{d}\bm{x} \\
													&=\sum_{T\in\mathscr{T}_i} \int_Tf(\bm{x})\phi_i(\bm{x})\ \mbox{d}\bm{x} \\
													&=\sum_{T\in\mathscr{T}_i} \int_Tf(\bm{x})p^{(T)}_{i^{(T)}}(\bm{x})\ \mbox{d}\bm{x} \\
													&=\sum_{T\in\mathscr{T}_i} \int_{\hat T}f(F^{(T)}(\bm{\hat x})\hat p_{i^{(T)}}(\bm{\hat x})|\mbox{det}(J)|\ \mbox{d}\bm{\hat x} \\
													&=h^2\sum_{T\in\mathscr{T}_i} \int_{\hat T}f(h\bm{\hat x}+v^{(T)}_1)\hat p_{i^{(T)}}(\bm{\hat x})\ \mbox{d}\bm{\hat x}						
\end{align*}

\section{Error Analysis}

In this section, we will discuss our expectations of the error and our method to compute the convergence order. In lecture, we had the following theorem\\
\\
\textbf{Theorem 4.29} \textit{Assume that the solution u of the given problem (5) satisfies
the regularity condition $u \in V=H^1(\Omega) \cap H^{k+1}(\Omega)$ with $k \geq m$. Then, The error in the solution of the finite element method satisfies the following bound:}
\[\|u-u_h\|_{m,\Omega}\leq ch^{k+1-m}|u|_{k+1,\Omega},\]
\textit{where c is a positive constant}\\
\\
The analytical solution is smooth, this enables us to chose $k$ freely. We chose $V$ to be a subspace of $H^1(\Omega)$, therefore $m=1$. The estimate in theorem 4.29 still contains a constant $c$ and can not be used directly. We will revise the computation of a convergence order from last semesters lecture:\\
Let $\alpha$ be the optimal exponent, such that equality holds
\begin{align}
\|u-u_h\|_{1,\Omega}&= ch^{\alpha}|u|_{k+1,\Omega} \\
\|u-u_{\frac{h}{2}}\|_{1,\Omega}&= c\frac{h^\alpha}{2^\alpha}|u|_{k+1,\Omega}
\end{align}
\\
Dividing (11) by (12) yields 

\[\frac{\|u-u_h\|_{1,\Omega}}{\|u-u_{\frac{h}{2}}\|_{1,\Omega}}= 2^\alpha\]
\[\Rightarrow  \alpha = \mbox{ln}\Big(\frac{\|u-u_h\|_{1,\Omega}}{\|u-u_{\frac{h}{2}}\|_{1,\Omega}}\Big)\cdot \mbox{ln}\Big(\frac{1}{2}\Big)\]
\end{document}
